\documentclass[11pt]{article}

\usepackage{booktabs,geometry,graphicx,hyperref,longtable,url}
\usepackage{subcaption,placeins}
\hypersetup{hidelinks}
\geometry{margin=1in}
\setlength{\parskip}{0.45em}
\setlength{\parindent}{0pt}

\title{Supplementary Material:\\Free-Boundary-Aware Risk Surfaces for American Options}
\author{SURF Computational Finance Project}
\date{\today}

\begin{document}
\maketitle

\section{Purpose}

This supplement supports the paper-facing manuscript by recording the
publication evidence audit, artifact locations, reproducibility commands,
diagnostics, and claim restrictions. It is intentionally more operational than
the main manuscript, but every reported result remains tied to an existing
project artifact.

\section{Publication Evidence Audit}

The Stage 15 publication evidence audit is generated from existing project artifacts by
\begin{verbatim}
python experiments/17_publication_evidence_audit.py
  --config configs/publication_evidence.yaml
\end{verbatim}
It writes the publication gap matrix, baseline comparison table, boundary
failure analysis, integrated scope summary, numerical benchmark summary, and
audit summary under
\path{results/experiment_17_publication_audit/}. This script reads
existing evidence and does not train models or rerun solver sweeps.

The audit classifies claims as main-text ready, supplementary, or limited.
Boundary and integrated-workflow claims are deliberately constrained because
boundary errors remain weak and the fresh integrated PSOR audit contains six
short-maturity cases rather than all 34 held-out test regimes.

\section{Figure Manifest}

Publication figures are generated by
\begin{verbatim}
python experiments/18_publication_figures.py
  --config configs/publication_figures.yaml
\end{verbatim}
The manifest is \texttt{results/figures/publication/figure\_manifest.csv}. Every figure used in the manuscript has both PNG and PDF outputs and a source artifact.

\section{Core Evidence Tables}

\begin{longtable}{p{0.28\linewidth}p{0.24\linewidth}p{0.38\linewidth}}
\toprule
Evidence area & Artifact & Manuscript use \\
\midrule
Solver validation & Stage 2 results & Main benchmark credibility \\
Grid/domain/runtime & Stage 4 results & Numerical limitation and stability evidence \\
Application sweep & Stage 5 results & Risk surfaces and delicate-regime audit \\
Dataset governance & Stage 6 results & Leakage-safe supervised labels \\
Baseline surrogates & Stage 7 results & Negative evidence for direct-price models \\
Positive premium & Stage 8 results & Obstacle-consistent price/premium result \\
Boundary heads & Stage 9 results & Scoped boundary diagnostic and limitation \\
Delta head & Stage 10 results & Strict-mask Delta diagnostic \\
Integrated workflow & Stage 11 results & Scenario-screening evidence only \\
\bottomrule
\end{longtable}

\section{Claim Restrictions}

The project supports scoped claims about audited CN/PSOR labels,
positive-premium obstacle consistency, supervised boundary diagnostics, bounded
Delta diagnostics, and scenario screening with PSOR audit. The workflow does
not replace CN/PSOR, and one neural network does not estimate every risk
quantity. The price model is not a Greek engine. A Gamma surrogate has not been
validated. Boundary weakness and failed regimes remain part of the evidence
package.

\section{Reproducibility Commands}

The following commands validate source integrity, regenerate the lightweight
publication audit and figures from existing summaries, and rerun the final claim
scan. They do not launch the expensive solver sweeps or neural training stages.
\begin{verbatim}
$env:PYTHONPATH="src;experiments"
$env:PYTHONPYCACHEPREFIX="$env:TEMP\american_option_pycache"

python -m unittest discover -s tests
python -m compileall -q src tests experiments
python experiments/17_publication_evidence_audit.py
  --config configs/publication_evidence.yaml
python experiments/18_publication_figures.py
  --config configs/publication_figures.yaml
python experiments/14_claim_audit.py --config configs/claim_audit.yaml
\end{verbatim}

Rounded values used in paper-facing prose are listed in
\texttt{docs/verified\_metric\_ledger.md}. Full-precision evidence remains in
the source CSV and JSON artifacts.

\section{Remaining Gaps}

The current package is a reviewer-facing computational research manuscript, not
a guarantee of journal acceptance. Stronger evidence would require a larger
fresh integrated PSOR audit, higher-resolution convergence checks, more reliable
boundary estimation, and external comparisons. Market calibration and
alternative dynamics would require new data provenance, mathematical
formulations, benchmark solvers, and validation gates; they are not completed
extensions of the present evidence.

% Auto-generated by experiments/19_report_visualization_enrichment.py
\section*{Visual Evidence Gallery}
\begin{figure}[!htbp]
\centering
\begin{subfigure}[t]{0.48\linewidth}
\centering
\includegraphics[width=\linewidth]{results/figures/report_enrichment/supplementary_material_full_gallery_workflow.pdf}
\caption{The supplement collects report visual evidence with a manifest for traceability.}
\end{subfigure}
\hfill
\begin{subfigure}[t]{0.48\linewidth}
\centering
\includegraphics[width=\linewidth]{results/figures/report_enrichment/supplementary_material_publication_workflow.pdf}
\caption{The publication workflow orients the supplement around CN/PSOR validation, clean labels, separate diagnostic heads, and continuing PSOR audit.}
\end{subfigure}
\end{figure}

\begin{figure}[!htbp]
\centering
\begin{subfigure}[t]{0.48\linewidth}
\centering
\includegraphics[width=\linewidth]{results/figures/report_enrichment/supplementary_material_publication_validation.pdf}
\caption{Solver validation evidence remains the foundation for all surrogate labels.}
\end{subfigure}
\hfill
\begin{subfigure}[t]{0.48\linewidth}
\centering
\includegraphics[width=\linewidth]{results/figures/report_enrichment/supplementary_material_publication_surrogate.pdf}
\caption{The surrogate comparison separates aggregate price error from payoff-obstacle consistency and does not treat the price model as a boundary or Greek estimator.}
\end{subfigure}
\end{figure}

\begin{figure}[!htbp]
\centering
\begin{subfigure}[t]{0.48\linewidth}
\centering
\includegraphics[width=\linewidth]{results/figures/report_enrichment/supplementary_material_publication_delta.pdf}
\caption{Delta evidence is separated from price model claims and evaluated under strict masks.}
\end{subfigure}
\hfill
\begin{subfigure}[t]{0.48\linewidth}
\centering
\includegraphics[width=\linewidth]{results/figures/report_enrichment/supplementary_material_publication_failure.pdf}
\caption{Failure and weak regime figures are retained to keep the final claim audit honest.}
\end{subfigure}
\end{figure}

\begin{figure}[!htbp]
\centering
\begin{subfigure}[t]{0.48\linewidth}
\centering
\includegraphics[width=\linewidth]{results/figures/report_enrichment/supplementary_material_supplement_matrix_pair.pdf}
\caption{The supplementary matrix pair gives readers a compact view of the A and B matrices used in the solver.}
\end{subfigure}
\hfill
\begin{subfigure}[t]{0.48\linewidth}
\centering
\includegraphics[width=\linewidth]{results/figures/report_enrichment/supplementary_material_supplement_lcp_vector.pdf}
\caption{The supplementary LCP diagram records the matrix inequality structure behind the obstacle diagnostics.}
\end{subfigure}
\end{figure}

\FloatBarrier


\end{document}
